\section{Secure Coding}
\label{sec:secure-coding}

El control A.14.2 de ISO 27001 establece que la organización debe establecer procesos de codificación segura para garantizar que el software sea desarrollado utilizando principios que reduzcan la introducción de vulnerabilidades durante la implementación \cite{iso27001}.

\subsection{Filosofía}

Muchos ataques exitosos no son sofisticados. Son errores simples como SQL Injection, XSS, Path Traversal, Command Injection o CSRF. Todos nacen del código. Por eso ISO 27001 crea un control exclusivo para la programación segura.

\subsection{Principios Fundamentales Aplicados a Workcodile-dev}

\subsubsection{1. Nunca confiar en la entrada}

La entrada del usuario siempre debe considerarse no confiable hasta que sea validada.

\textbf{Ejemplo en Workcodile-dev (backend -- Express):}

\begin{lstlisting}[language=Java, caption={Código inseguro -- sin validación}]
// INSEGURO - No validación de entrada
app.get('/api/posts/:id', (req, res) => {
  const post = Post.findById(req.params.id);
  res.json(post);
});
\end{lstlisting}

\textbf{Solución implementada:}

\begin{lstlisting}[language=Java, caption={Código seguro -- con validación de ObjectId}]
// SEGURO - Con validación de ObjectId
app.get('/api/posts/:id', (req, res) => {
  if (!mongoose.Types.ObjectId.isValid(req.params.id)) {
    return res.status(400).json({ error: 'ID invalido' });
  }
  const post = Post.findById(req.params.id);
  res.json(post);
});
\end{lstlisting}

\subsubsection{2. Validar siempre}

Validar significa comprobar que el dato cumple lo esperado antes de procesarlo. Se validan: tipo, longitud, formato, rango y conjunto permitido de valores.

\textbf{Ejemplo en Workcodile-dev (registro de usuario):}

\begin{lstlisting}[language=Java, caption={Middleware de validación con express-validator}]
const { body, validationResult } = require('express-validator');

app.post('/api/auth/register', [
  body('email').isEmail().normalizeEmail(),
  body('password').isLength({ min: 8 })
    .matches(/^(?=.*[A-Z])/),
  body('name').trim().escape()
    .isLength({ max: 100 }),
], (req, res) => {
  const errors = validationResult(req);
  if (!errors.isEmpty()) {
    return res.status(400).json({
      errors: errors.array()
    });
  }
});
\end{lstlisting}

\subsubsection{3. Sanitizar cuando corresponda}

Sanitizar significa preparar los datos para su uso seguro. Un input como \texttt{<script>alert(1)</script>} puede producir Stored XSS si se almacena sin control.

\textbf{En Workcodile-dev (frontend -- React):}

\begin{lstlisting}[language=Java, caption={Sanitización en React}]
// SEGURO - React escapa automáticamente
const UserBio = ({ bio }) => {
  return <div>{bio}</div>;
};

// INSEGURO - Nunca usar sin sanitizar
const UserBio = ({ bio }) => {
  return <div
    dangerouslySetInnerHTML={{__html: bio}}
  />;
};
\end{lstlisting}

\subsubsection{4. Principio de mínimo privilegio}

El software debe usar únicamente los permisos que necesita para operar.

\textbf{En Workcodile-dev:}
\begin{itemize}
    \item MongoDB: usuario con permisos solo sobre la BD de la aplicación.
    \item MinIO: credenciales con acceso limitado al bucket de uploads.
    \item Docker: usuario no root, puertos mínimos expuestos.
\end{itemize}

\subsubsection{5. Gestión de secretos}

Nunca almacenar secretos en código fuente ni en repositorios.

\textbf{En Workcodile-dev:}
\begin{itemize}
    \item Archivo \texttt{.env} con credenciales (excluido de Git).
    \item Archivo \texttt{.env.example} como plantilla versionada.
    \item Variables de entorno cargadas con \texttt{dotenv}.
    \item Nunca: \texttt{const PW = "admin123"} en código.
\end{itemize}

\subsubsection{6. Manejo seguro de errores}

Los mensajes de error no deben exponer información sensible del sistema.

\begin{lstlisting}[language=Java, caption={Manejo seguro de errores}]
// INSEGURO - Expone detalles del error
app.use((err, req, res, next) => {
  res.status(500).json({
    error: err.message,
    stack: err.stack  // NUNCA exponer
  });
});

// SEGURO - Mensaje generico, log interno
app.use((err, req, res, next) => {
  console.error(`[${new Date()}] ${err.message}`);
  res.status(500).json({
    error: 'Error interno del servidor'
  });
});
\end{lstlisting}

\subsubsection{7. Dependencias seguras}

Gran parte de los ataques actuales no vienen del código propio, sino de dependencias externas.

\textbf{En Workcodile-dev:}
\begin{itemize}
    \item Ejecutar \texttt{npm audit} periódicamente.
    \item Actualizar paquetes con vulnerabilidades conocidas.
    \item Eliminar dependencias abandonadas.
    \item Usar \texttt{npm audit fix} para correcciones automáticas.
\end{itemize}

\subsection{OWASP Top 10 y Workcodile-dev}

El control A.14.2 gira alrededor de prevenir las vulnerabilidades del OWASP Top 10 \cite{owasp2021}:

\begin{table}[h]
\centering
\small
\resizebox{\columnwidth}{!}{%
\begin{tabular}{|l|l|p{5cm}|}
\hline
\textbf{OWASP} & \textbf{¿Aplica?} & \textbf{Controles} \\
\hline
Injection & Sí & Validación ObjectId, consultas parametrizadas \\
Broken Access Control & Sí & Middleware JWT por roles \\
Cryptographic Failures & Sí & bcrypt, JWT con expiración \\
Security Misconfiguration & Sí & .env, Docker hardening \\
Vulnerable Components & Sí & npm audit, actualización deps \\
Authentication Failures & Sí & JWT expiración, bcrypt cost \\
\hline
\end{tabular}%
}
\caption{Mapeo OWASP Top 10 a Workcodile-dev}
\label{tab:owasp}
\end{table}
